\chapter{Planteamiento del Problema y Justificación}\label{cap:problema}
Este capítulo delimita el problema que motiva el proyecto, formula la pregunta de investigación y expone las razones académicas, técnicas y sociales que justifican su desarrollo.

\section{Planteamiento del problema}
El aprendizaje automático suele enseñarse y practicarse mediante \emph{frameworks} de alto nivel que ocultan los detalles de su funcionamiento interno. Con unas pocas líneas de código es posible entrenar un modelo, pero sin comprender cómo se representan los datos en vectores y matrices, cómo se calcula la salida de cada neurona, cómo se mide el error ni cómo se ajustan los pesos. Esta situación genera profesionales que dependen de herramientas cuyo comportamiento no pueden explicar, depurar ni optimizar.

En el contexto local, los estudiantes de los primeros semestres de la FICCT carecen de materiales didácticos que conecten de forma explícita los contenidos de Programación 1 y del álgebra lineal con el funcionamiento real de una red neuronal. Además, buena parte de los recursos disponibles se basa en lenguajes interpretados y bibliotecas externas, lo que dificulta observar el costo computacional real de cada operación.

\subsection{Pregunta de investigación}
¿Es posible construir, con fines educativos y utilizando únicamente la biblioteca estándar del lenguaje Rust, una red neuronal artificial multicapa capaz de aprender relaciones lineales, resolver un problema no lineal como XOR y clasificar patrones visuales simples como las vocales A, E, I, O, U, verificando su correcto funcionamiento mediante pruebas automatizadas y experimentos reproducibles?

\subsection{Delimitación}
\begin{itemize}
\item \textbf{Alcance temático:} redes neuronales densas (\emph{feedforward}) con funciones de activación Sigmoid, ReLU y lineal, pérdida MSE y descenso de gradiente estocástico.
\item \textbf{Alcance técnico:} implementación en Rust (edición 2024) sin crates externos; ejecución en un computador personal convencional.
\item \textbf{Alcance experimental:} tres experimentos progresivos ($y=2x$, XOR y vocales) con datos sintéticos definidos explícitamente en el código.
\item \textbf{Fuera de alcance:} redes convolucionales o recurrentes, aceleración por GPU y conjuntos de datos masivos.
\end{itemize}

\section{Justificación}
\subsection{Justificación académica}
Construir una red neuronal desde cero obliga a traducir cada ecuación matemática en código verificable. Ese ejercicio consolida conocimientos de álgebra lineal (productos matriz-vector, transpuestas, producto de Hadamard), de cálculo diferencial (regla de la cadena) y de programación estructurada, y permite comprender por qué funcionan técnicas que en los \emph{frameworks} aparecen como una sola llamada a función \citep{nielsen2015,deisenroth2020}.

\subsection{Justificación técnica}
Rust ofrece rendimiento comparable al de C sin recolector de basura y con seguridad de memoria garantizada por el compilador \citep{klabnik2023,jung2018}. Su sistema de tipos y el manejo explícito de errores mediante \texttt{Result} permiten detectar en tiempo de ejecución, de forma controlada, errores de dimensión que en otros lenguajes producirían resultados incorrectos silenciosos. Además, su herramienta \texttt{cargo} integra compilación, pruebas y documentación, lo que facilita una metodología de desarrollo verificable.

\subsection{Justificación social y educativa}
El proyecto se publica como código abierto y no requiere hardware especializado, licencias ni conexión permanente a servicios en la nube. Esto lo vuelve accesible para estudiantes de cualquier provincia y lo convierte en un material reutilizable para la enseñanza de la IA en la universidad pública.

\subsection{Viabilidad}
El proyecto es viable porque utiliza herramientas gratuitas (Rust, Cargo, Git, \LaTeX), se ejecuta en equipos convencionales y se apoya en una bibliografía consolidada. La validación mediante pruebas automatizadas y experimentos reproducibles con semilla fija garantiza que los resultados puedan ser verificados por cualquier evaluador.
